\documentclass[10pt,a4paper]{amsart}
\usepackage[margin=19mm]{geometry}
\usepackage{amsmath,amssymb,mathtools}
\usepackage[scr=rsfs,cal=euler]{mathalfa}
\usepackage{xfrac}
\usepackage[unicode,hidelinks,bookmarks=false]{hyperref}
\newcommand{\ie}{i.e.\ }
\newcommand{\cf}{cf.\ }
\newcommand{\resp}{respectively\ }
\newcommand{\Subsection}{\subsection}
\providecommand{\nobreakdash}{\nobreak\hbox{-}}
\setlength{\emergencystretch}{2em}
\multlinegap0pt
\usepackage{bm}
\usepackage{bbm}
\usepackage{dsfont}
\usepackage{xspace}
\usepackage{cancel}
\usepackage{mathtools}
\usepackage{cleveref}
\usepackage{amsthm}
\makeatletter
\@ifundefined{lemma}{\newtheorem{lemma}{Lemma}}{}
\@ifundefined{theorem}{\newtheorem{theorem}{Theorem}}{}
\makeatother
\NewDocumentCommand{\CoeffSpace}{o}{\mathfrak{B}\IfValueT{#1}{^{(#1)}}}
\NewDocumentCommand{\RadHaarSpace}{o}{\IfValueTF{#1}{\mathcal{G}_{#1}}{\mathcal{G}}}
\NewDocumentCommand{\card}{m}{\ensuremath{\operatorname{card}\left(#1\right)}}
\NewDocumentCommand{\coeffC}{s o o}{
    \IfValueTF{#3}
        {\IfBooleanTF{#1}{\bar{c}}{c}_{#3}}
        {\bm{\IfBooleanTF{#1}{\bar{c}}{c}}}%
    \IfValueT{#2}{^{(#2)}}
}
\NewDocumentCommand{\coeff}{s o o}{
    \IfValueTF{#3}
        {\IfBooleanTF{#1}{\bar{b}}{b}_{#3}}
        {\bm{\IfBooleanTF{#1}{\bar{b}}{b}}}%
    \IfValueT{#2}{^{(#2)}}
}
\NewDocumentCommand{\decoder}{d<> o}{\ensuremath{\delta_{D\IfValueT{#1}{,#1}}\IfValueT{#2}{^{#2}}}}
\NewDocumentCommand{\dyadicU}{m m}{\set{U}_{#2}^{(#1)}}
\NewDocumentCommand{\encoder}{d<> o}{\ensuremath{\delta_{E\IfValueT{#1}{,#1}}\IfValueT{#2}{^{#2}}}}
\NewDocumentCommand{\ld}{}{\operatorname{ld}}
\NewDocumentCommand{\nn}{o}{\ensuremath{\mathbb{N}\IfValueT{#1}{^{#1}}}}
\NewDocumentCommand{\norm}{s m o}{
    \ensuremath{
        \IfBooleanTF{#1}{\|{#2}\|}{\left\|{#2}\right\|}
        \IfValueT{#3}{_{#3}}
    }
}
\NewDocumentCommand{\rr}{o}{\ensuremath{\mathbb{R}\IfValueT{#1}{^{#1}}}}
\NewDocumentCommand{\sepClass}{o}{\mathscr{C}_{\mathrm{sep}}\IfValueT{#1}{^{#1}}}
\NewDocumentCommand{\walsh}{s m m d()}{\ensuremath{\IfBooleanTF{#1}{\bar{\omega}}{\omega}_{#3}^{(#2)}\IfValueT{#4}{\!\left(#4\right)}}}
\NewDocumentCommand{\wavelet}{o o d()}{\ensuremath{\varphi\IfValueT{#2}{_{#2}}\IfValueT{#1}{^{(#1)}}\IfValueT{#3}{\!\left(#3\right)}}}
\newcommand{\defeq}{\mathbin{:=}}
\newcommand{\set}[1]{\ensuremath{\mathcal{#1}}}
\title[The Information-Theoretic Benefit of Shared Representations ]{The Information-Theoretic Benefit of Shared Representations under Orthogonality Constraints}
\author{Thomas Dittrich, Oliver Potocki, Philipp Grohs}
\date{Source: arXiv:2606.16028v1 (CC BY 4.0)}
\begin{document}
\maketitle

\noindent\emph{Source and licence.} The Information-Theoretic Benefit of Shared Representations under Orthogonality Constraints. Version arXiv:2606.16028v1 (CC BY 4.0), distributed under the Creative Commons Attribution 4.0 International licence, as recorded in the arXiv licence metadata for this exact version and shown on the abstract page https://arxiv.org/abs/2606.16028v1. Full rights notice: licenses/arxiv-2606.16028-CC-BY-4.0.txt.\par\smallskip

\noindent\textit{Editorial scope.} Counted unit: the Theorem whose source label is \texttt{thm:wavelet\_separated} in the article (source0.tex lines 965--979, source label \texttt{thm:wavelet\_separated}), delivered together with the article's own complete original proof of it (source0.tex lines 981--1124, one \texttt{proof} environment of 144 lines). No mathematics is written, repaired or re-derived: statement and proof are the article's own text line for line. Required context retained uncounted: thm:wavelet\_distance\_lower (lines 596--606). Every definition, display and label of those lines is retained unchanged. Omitted from this package and not redistributed: 1--595, 607--964, 980--980, 1125--1504. Nothing inside a retained excerpt is omitted or altered: every retained body is delivered exactly as the article's lines are written, so no omission receipt of any kind accompanies this package. Citations: the retained excerpts carry no citation command at all, so no bibliography entry of the article is transcribed and no reference is redistributed. Reference adaptation: none was needed and none was made; every reference inside the delivered unit resolves inside it, so no unresolved reference or citation is produced. Printed slips kept literal: no printed word, symbol, hypothesis or number of the counted statement or of its proof was altered. Layout and class: the article's own class is \texttt{article} with packages including jmlr2e, amsmath, bm, mathrsfs, cleveref, lastpage; the wrapper is the lane's pinned \texttt{amsart} editorial wrapper at 10pt on A4, which restates the theorem-like environments the retained text needs and adds the article's own macro definitions that the retained text needs (\texttt{\textbackslash CoeffSpace}, \texttt{\textbackslash RadHaarSpace}, \texttt{\textbackslash card}, \texttt{\textbackslash coeff}, \texttt{\textbackslash coeffC}, \texttt{\textbackslash decoder}, \texttt{\textbackslash defeq}, \texttt{\textbackslash dyadicU}, \texttt{\textbackslash encoder}, \texttt{\textbackslash ld}, \texttt{\textbackslash nn}, \texttt{\textbackslash norm}, \texttt{\textbackslash rr}, \texttt{\textbackslash sepClass}, \texttt{\textbackslash set}, \texttt{\textbackslash walsh}, \texttt{\textbackslash wavelet}), with extra wrapper packages (bm, bbm, dsfont, xspace, cancel, mathtools, cleveref). The delivered numbering is the unit's own. Assets: no retained span contains a figure environment, a graphics-inclusion command, a PDF-page inclusion, a table or a TikZ picture; no image file is copied into this package, the package declares no asset dependency and no image file is redistributed with it. Licence: version arXiv:2606.16028v1 of this article is distributed under the Creative Commons Attribution 4.0 International licence, as recorded in the arXiv licence metadata for this exact version and shown on the abstract page https://arxiv.org/abs/2606.16028v1. A full rights notice is delivered at licenses/arxiv-2606.16028-CC-BY-4.0.txt. Characters: every retained excerpt is pure ASCII, so the delivered engine, XeLaTeX with its default Latin Modern text font, reports no missing character.
\begin{lemma}[Minimum Distance of Wavelet Series]
    \label[lemma]{thm:wavelet_distance_lower}
    Let $L_0\in\nn_0$, $j\in[2^{L_0}]$, and
    $\coeff,\coeffC\in\CoeffSpace$ be arbitrary but fixed
    and set\footnote{Note that we use the convention $\min \emptyset=-\max \emptyset=\infty$.}
    $L\defeq\min\{l\in\nn_0|\exists x\in\dyadicU{L_0}{j}:g_l(x;\coeff)\neq g_l(x;\coeffC)\}$,
    then
    \begin{align}
        \norm{g(\cdot;\coeff)-g(\cdot;\coeffC)}[L^\infty\!\left(\dyadicU{L_0}{j}\right)]\geq 2^{-\max\{L,L_0-1\}}.
    \end{align}
\end{lemma}

\begin{theorem}[Complexity of Separate Approximation]
    \label{thm:wavelet_separated}
    Let $\set{G}$ be the class of Rademacher-Haar wavelets, $M\in\nn$, $n\in\nn[M]$, $L=\lceil\ld(M)\rceil$, $(f_0,\ldots,f_{M-1})$ be the first $M$ Walsh-Sawtooth functions of order $L$, and $\set{F}\defeq \{(f_0\circ g,\ldots,f_{M-1}\circ g)|g\in\set{G}\}$.
    Then, with $N=\min_{i\in[M]}n_i$,
    \begin{align}
        \inf_{(\encoder,\decoder)\in\sepClass[n]}
        \sup_{h\in\set{F}}
        \norm{h-\decoder\circ\encoder(h)}[{L^\infty([0,1];\rr[M])}]
        \geq
        \sqrt{3}
        2^{L-\lceil\ld(N+2^{L})\rceil}
        \geq
        \frac{2^{L-1}\sqrt{3}}{N+2^{L}}.
    \end{align}
\end{theorem}

\begin{proof}
In separated approximation, we discretize each element of the $M$-tuple $F\in\set{F}$ individually.
Instead of considering the class $\sepClass[n]$ of encoder/decoder pairs directly for the full tuple, we first consider element-wise discretization maps $\Delta_{\set{C}}$ with $\set{C}\subset L^\infty([0,1])$ and then combine the element-wise distortion, measured in $L^\infty([0,1])$, into the distortion of the full tuple, measured in $L^\infty([0,1];\rr[M])$.
We then have the equivalent formulation
\begin{align}
    &\kern-2em\inf_{(\encoder,\decoder)\in\sepClass[n]}
    \sup_{g\in\set{G}}
    \norm{g-\decoder\circ\encoder(g)}[{L^\infty([0,1];\rr[M])}]
    \\&
    =
    \inf_{\substack{\set{C}_i\subset L^\infty([0,1]),\\\card{\set{C}}=2^{n_i},\\i\in[M]}}
    \sup_{h\in\set{F}}
    \max_{j\in[M]}
    \norm{h_j-\Delta_{\set{C}_j}(h_j)}[{L^\infty([0,1])}]
    \\&
    =
    \inf_{\substack{\set{C}_i\subset L^\infty([0,1]),\\\card{\set{C}}=2^{n_i},\\i\in[M]}}
    \max_{j\in[M]}
    \sup_{g\in\set{G}}
    \norm{f_j\circ g-\Delta_{\set{C}_j}(f_j\circ g)}[{L^\infty([0,1])}]
    \\&
    =
    \max_{j\in[M]}
    \inf_{\substack{\set{C}_j\subset L^\infty([0,1]),\\\card{\set{C}}=2^{n_j}}}
    \sup_{g\in\set{G}}
    \norm{f_j\circ g-\Delta_{\set{C}_j}(f_j\circ g)}[{L^\infty([0,1])}].
\end{align}
Here, we first simply swap the supremum and the maximum (that is, two suprema) and simplify the expression.
Second, we see that the infimum is over a Cartesian product of families of sets and the argument of the maximum only depends on a single element in the Cartesian product.
Therefore, we can also swap the infimum and the maximum.

We now continue with the analysis for a single index $i\in[M]$.
Let $\coeff\in\CoeffSpace$ be fixed, and choose $\coeffC\in\CoeffSpace$ such that
\begin{alignat}{2}
    \coeffC[l][j]&=\coeff[l][j]\cdot\walsh{L}{i}(g_{L-1}(j2^{-l};\coeff))\qquad&&\text{for}\qquad l\geq L,\,j\in[2^l].
\end{alignat}
Then for all $x\in [0,1)$ we get
\begin{align}
    (2^{L}\sqrt{3})^{-1}\cdot f_i\circ g(x;\coeff)&=(g(x;\coeff)-g_{L-1}(x;\coeff))(\omega_i\circ g_{L-1}(x;\coeff))
    \\&
    =
    \sum_{l=L}^\infty\sum_{j=0}^{2^l-1}\coeff[l][j]\wavelet[l][j](x) \walsh{L}{i}(g_{L-1}(x;\coeff))
    \\&
    =
    \sum_{l=L}^\infty\sum_{j=0}^{2^l-1}\left(\coeff[l][j] \walsh{L}{i}(g_{L-1}(j2^{-l};\coeff))\right)\wavelet[l][j](x)
    \\&
    =
    \sum_{l=L}^\infty\sum_{j=0}^{2^l-1}\coeffC[l][j]\wavelet[l][j](x)
    \\&
    =
    g(x;\coeffC)-g_{L-1}(x;\coeffC).
\end{align}
Here we used that the wavelets $\wavelet[l][j]$ are only supported on $\dyadicU{l}{j}$ and $g_{L-1}$ is piecewise constant such that $g_{L-1}(x;\coeff)=g_{L-1}(j2^{-l};\coeff)$ holds for all $x\in\dyadicU{l}{j}$.
As such an equivalent construction exists for every element of $\CoeffSpace$, the discretization with $B\in\nn$ bits reduces to
\begin{align}
    &\kern-2em\inf_{\substack{\set{C}\subset L^\infty(\dyadicU{0}{0})\\\card{\set{C}}=2^B}}\sup_{g\in\RadHaarSpace}\norm{f_i\circ g-\Delta_{\set{C}}(f_i\circ g)}[{L^\infty(\dyadicU{0}{0})}]
    \\&
    =
    2^L\sqrt{3}\inf_{\substack{\set{C}\subset L^\infty(\dyadicU{0}{0})\\\card{\set{C}}=2^B}}\sup_{\coeffC\in\CoeffSpace}\norm{g(\cdot;\coeffC)-g_{L-1}(\cdot;\coeffC)-\Delta_{\set{C}}(g(\cdot;\coeffC)-g_{L-1}(\cdot;\coeffC))}[{L^\infty(\dyadicU{0}{0})}].
\end{align}

We now show for an arbitrary $L_1\geq L$ that if $\set{C}$ forms an $\varepsilon$-cover of $\{g(\cdot;\coeff)-g_{L-1}(\cdot;\coeff)|\coeff\in\CoeffSpace\}$ with $\varepsilon<2^{-L_1-1}$ then $\ld(\card{C})> 2^{L_1+2}-2^{L+1}$.

Let
\begin{align}
    \CoeffSpace[L:L_1]
    &\defeq
    \{\coeff\in\CoeffSpace[L_1]|\forall l<L:\coeff[l]=1\},
    \\
    N
    &\defeq
    \card{\CoeffSpace[L:L_1]}=2^{\sum_{l=L}^{L_1}2^l}=2^{2^{L_1+1}-2^L},
\end{align}
and consider an enumeration
\begin{align}
    \CoeffSpace[L:L_1]=\{\coeff*_k|k\in[N]\}.
\end{align}
We then construct $\coeff_k,\coeffC_k\in\CoeffSpace$ for $k\in[N]$ such that
\begin{alignat}{2}
    \coeff[l][k,j]&=\coeffC[l][k,j]=\coeff*[l][k,j]&\qquad&\text{for}\qquad l\leq L_1,\, j\in[2^l],\\
    \coeff[l][k,j]&=-\coeffC[l][k,j]=1&&\text{for}\qquad l>L_1,\,j\in[2^l].
\end{alignat}
We see for $k_1,k_2\in[N]$ with $k_1\neq k_2$ that there exist $x\in\dyadicU{0}{0}$ and $L_2\in[L_1+1]$ with $g_{L_2}(x;\coeff_{k_1})=g_{L_2}(x;\coeffC_{k_1})\neq g_{L_2}(x;\coeff_{k_2})=g_{L_2}(x;\coeffC_{k_2})$.
From \cref{thm:wavelet_distance_lower} we then get
\begin{align}
    \norm{g(\cdot;\coeff_{k_1})-g(\cdot;\coeff_{k_2})}[L^\infty(\dyadicU{0}{0})]
    &\geq 2^{-L_2}\geq 2^{-L_1}
    \\
    \norm{g(\cdot;\coeff_{k_1})-g(\cdot;\coeffC_{k_2})}[L^\infty(\dyadicU{0}{0})]
    &\geq 2^{-L_2}\geq 2^{-L_1}
    \\
    \norm{g(\cdot;\coeffC_{k_1})-g(\cdot;\coeff_{k_2})}[L^\infty(\dyadicU{0}{0})]
    &\geq 2^{-L_2}\geq 2^{-L_1}
    \\
    \norm{g(\cdot;\coeffC_{k_1})-g(\cdot;\coeffC_{k_2})}[L^\infty(\dyadicU{0}{0})]
    &\geq 2^{-L_2}\geq 2^{-L_1}.
\end{align}
Furthermore, for any $k\in[N]$ we get
\begin{align}
    g(0;\coeff_k)-g(0;\coeffC_k)
    =
    \sum_{l=L_1+1}^\infty 2^{-l-1}(1-(-1))
    =
    \sum_{l=L_1+1}^\infty 2^{-l}
    =2^{-L_1}.
\end{align}
For an $\varepsilon$-cover with $\varepsilon<2^{-L_1-1}$ and center points $\set{C}$ we get that no pair of elements in
\begin{align}
    \{\coeff_k,\coeffC_k|k\in[N]\}
\end{align}
can be in the same ball of radius $\varepsilon$ and therefore $\card{\set{C}}\geq 2N>N=2^{2^{L_1+1}-2^{L}}$.
From the reverse implication we see for an arbitrary $B\in\nn$ with $B\leq 2^{L_1+1}-2^{L}$ for some $L_1\geq L$ that 
\begin{align}
    \inf_{\substack{\set{C}\subset L^\infty(\dyadicU{0}{0})\\\card{\set{C}}=2^B}}\sup_{g\in\RadHaarSpace}\norm{f_i\circ g-\Delta_{\set{C}}(f_i\circ g)}[{L^\infty(\dyadicU{0}{0})}]\geq \sqrt{3}2^{L-L_1-1}.
\end{align}
Specifically, for any $B\in\nn$ we can choose $\overline{L}_1\defeq\lceil\ld(B+2^{L})\rceil-1\leq \ld(B+2^{L})$ and get $B\leq 2^{\overline{L}_1+1}-2^{L}$, thus,
\begin{align}
    &\kern-2em\inf_{\substack{\set{C}\subset L^\infty(\dyadicU{0}{0})\\\card{\set{C}}=2^B}}\sup_{g\in\RadHaarSpace}\norm{f_i\circ g-\Delta_{\set{C}}(f_i\circ g)}[{L^\infty(\dyadicU{0}{0})}]
    \\&
    \geq
    \sqrt{3}
    2^{L-\overline{L}_1-1}
    \geq
    \sqrt{3}2^{L-\ld(B+2^{L})-1}
    \geq
    \sqrt{3}\frac{2^{L-1}}{B+2^{L}}
\end{align}

Overall, with $N=\min_{i\in[M]}n_i$, 
we get
\begin{align}
    &\kern-2em\inf_{(\encoder,\decoder)\in\sepClass[n]}
    \sup_{h\in\set{F}}
    \norm{h-\decoder\circ\encoder(h)}[{L^\infty([0,1];\rr[M])}]
    \\&
    =
    \max_{j\in[M]}\inf_{\substack{\set{C}\subset L^\infty([0,1])\\\card{\set{C}}=2^{n_j}}}\sup_{g\in\RadHaarSpace}\norm{f_j\circ g-\Delta_{\set{C}}(f_j\circ g)}[{L^\infty([0,1])}]
    \\&
    \geq
    \sqrt{3}2^{L-\lceil\ld(N+2^{L})\rceil}
    \geq 
    \frac{2^{L-1}\sqrt{3}}{N+2^{L}}.
\end{align}
\end{proof}

\end{document}
